% !TeX root = ../informe.tex
% ===========================================================================
%  3. Alcance funcional
%
%  Qué hace el sistema y qué deliberadamente no hace. Es contexto para las
%  secciones de arquitectura, no el capítulo principal: el documento de
%  alcance del curso ya describe estas funcionalidades, así que aquí basta
%  con recogerlas de forma compacta y señalar lo que el equipo acotó.
% ===========================================================================

\section{Alcance funcional}
\label{sec:alcance}

\subsection{Roles y permisos}
\label{sec:roles}

\instruccion{Una matriz de funcionalidad por rol se lee de un vistazo y
evita párrafos repetitivos. Marca con «Sí» y «No»; si algún permiso depende
de una condición (cancelar solo las propias), dilo en la celda.}

\begin{table}[H]
  \centering
  \caption{Matriz de roles y permisos.}
  \label{tab:roles}
  \begin{tabular}{L{6.2cm} C{2.6cm} C{2.6cm}}
    \toprule
    \textbf{Funcionalidad} & \textbf{Usuario final} & \textbf{Administrador} \\
    \midrule
    Registro e inicio de sesión                     & Sí & Sí \\
    Consultar disponibilidad de canchas             & Sí & Sí \\
    Crear una reserva                               & Sí & No \\
    Cancelar una reserva                            & Solo las propias & Cualquiera \\
    Consultar el historial de reservas propias      & Sí & No \\
    Consultar el listado global de reservas         & No & Sí \\
    Gestionar el catálogo de canchas                & No & Sí \\
    Gestionar horarios y bloqueos de mantenimiento  & No & Sí \\
    Gestionar usuarios (activar / inactivar)        & No & Sí \\
    Visualizar los reportes de ocupación            & No & Sí \\
    \bottomrule
  \end{tabular}
\end{table}

\subsection{Módulos y pantallas}
\label{sec:modulos}

\instruccion{Los módulos funcionales con sus pantallas. Conviene que esta
tabla case con los microfrontends de la \cref{tab:microfrontends}: si un
módulo no corresponde a ningún remote, o al revés, es señal de que la
descomposición no está clara.}

\begin{table}[H]
  \centering
  \caption{Módulos y pantallas del sistema.}
  \label{tab:modulos}
  \begin{tabular}{L{3.2cm} L{4.4cm} L{5.4cm}}
    \toprule
    \textbf{Módulo} & \textbf{Pantalla} & \textbf{Descripción} \\
    \midrule
    Seguridad      & Inicio de sesión / Registro & Autenticación de ambos roles y alta de cuentas nuevas. \\
    Reservas       & Consulta de disponibilidad  & Rejilla de bloques por cancha y fecha, libres y ocupados. \\
    Reservas       & Nueva reserva               & Confirmación del bloque elegido. \\
    Reservas       & Mis reservas                & Listado propio, con la opción de cancelar. \\
    Administración & Gestión de canchas          & Alta, edición, horario, bloqueos y estado de cada cancha. \\
    Administración & Gestión de reservas         & Listado global, con la opción de cancelar cualquiera. \\
    Administración & Gestión de usuarios         & Activar e inactivar cuentas. \\
    Reportes       & Reportes y estadísticas     & Los cuatro indicadores de ocupación y demanda. \\
    \bottomrule
  \end{tabular}
\end{table}

\subsection{Fuera de alcance}
\label{sec:fuera-alcance}

\instruccion{Enumera lo excluido. No es una disculpa: delimitar el alcance
es una decisión de diseño, y declararla evita que se lea como una carencia.}

Queda fuera del alcance, por decisión explícita y no por olvido:

\begin{itemize}
  \item \textbf{Pasarela de pagos o cobro en línea.} El club cobra por sus
        medios; la reserva no es una transacción económica en este sistema.
  \item \textbf{Notificaciones por correo, SMS o push.} Las confirmaciones y
        los rechazos se comunican \emph{en pantalla}, en el momento. Añadir un
        canal asíncrono traería una infraestructura —cola, reintentos,
        plantillas— que ningún criterio de la rúbrica evalúa.
  \item \textbf{Reservas recurrentes.} Cada reserva es un bloque concreto en
        una fecha concreta.
  \item \textbf{Torneos, ligas e inscripciones grupales.} La unidad reservable
        es un bloque para un usuario, no un evento con participantes.
  \item \textbf{Aplicación móvil nativa.} El sistema es web y responsive; se
        usa desde el navegador del teléfono.
  \item \textbf{Reportes analíticos o de inteligencia de negocio.} Los cuatro
        indicadores de §\ref{sec:alcance} se calculan sobre los datos vivos y se
        muestran en pantalla; no hay exportación a formatos analíticos.
\end{itemize}

Esta lista importa tanto como la de funcionalidades incluidas. Delimitar el
alcance es lo que permite que lo que sí está construido esté \emph{terminado}:
el tiempo que se gastara en cualquiera de estos seis puntos saldría del que
sostiene las ocho reglas de negocio y sus pruebas.
